%% kapa: exam=1 round=33 year=2022 period=1 form=A subject=economics subject_ko=경제학원론 seq=13 number=53
%% source: Q-Net 감정평가사 기출문제 (cst00309) artlSeq=5209997 file=raw/1008/5209997/2022년도 제33회 감정평가사 1차시험 1교시(민법,경제학,부동산학).pdf
%% conversion: kordoc 4.24.51 -> tools/gen1.py (auto); review=proofread-sonnet; flags=none
\begin{question}{53}
X재 가격이 하락할 때 아래의 설명 중 옳은 것을 모두 고른 것은? (단, X재와 Y재만 존재하며 주어진 소득을 두 재화에 모두 소비한다.)
\begin{framedtext}
ㄱ. X재가 정상재인 경우 보상수요곡선은 보통수요곡선보다 더 가파르게 우하향하는 기울기를 가진다. \par
ㄴ. X재가 열등재인 경우 보상수요곡선은 우상향한다. \par
ㄷ. X재가 기펜재인 경우 보통수요곡선은 우상향하고 보상수요곡선은 우하향한다.
\end{framedtext}
\begin{choices}
  \item ㄱ
  \item ㄴ
  \item ㄱ, ㄷ
  \item ㄴ, ㄷ
  \item ㄱ, ㄴ, ㄷ
\end{choices}
\end{question}
\answer{3}
